% 設定頁面邊距
\geometry{
    a4paper,
    left=1cm,
    right=1cm,
    top=2cm,
    bottom=2cm
}

\hypersetup{hidelinks}
\setlength{\headheight}{14.5pt}

% 允許 TeX 在必要時提早折行（吸收少量過長行），避免文字溢出右邊界。
\emergencystretch=3em

% 設定字型
\setmainfont[
    Path = ./fonts/Times_New_Roman/,
    BoldFont = *-bold,
    ItalicFont = *-italic,
    BoldItalicFont = *-bold_italic
]{times}

\setmonofont[
    Path = ./fonts/Roboto_Mono/static/,
    UprightFont = *-Regular,
    BoldFont = *-Bold,
    ItalicFont = *-Italic,
    BoldItalicFont = *-BoldItalic
]{RobotoMono}

\xeCJKsetup{AutoFakeBold=true, AutoFakeSlant=true}
\setCJKmainfont[
    Path = ./fonts/kaiu/,
]{kaiu}
\setCJKmonofont{./fonts/kaiu/kaiu.ttf}

% 頁面設定
\fancyhf{}
\rfoot{\thepage}

% 設置圖片路徑
\graphicspath{{fig/}}

% 標題樣式
\titleformat{\section}{\large\bfseries}{\thesection}{1em}{}
\titleformat{\subsection}{\normalsize\bfseries}{\thesubsection}{1em}{}


\pagestyle{fancy}
\fancyhf{}
\rhead{June. 2026}
\lhead{\textbf{Visual Recognition using Deep Learning
2026 Spring}}
\rfoot{\thepage}
